\section{Pendahuluan}

\subsection{Latar Belakang}
Sistem tertanam dapat digunakan untuk memantau kondisi lingkungan melalui sensor dan mengendalikan komponen berdasarkan hasil pembacaan tersebut. Salah satu penerapannya adalah sistem kontrol dan deteksi udara serta gas yang memanfaatkan sensor untuk mendeteksi kondisi atau kandungan tertentu di lingkungan.

Dalam FGD ini, pembahasan difokuskan pada eksplorasi sensor, aktuator, dan komponen pendukung yang digunakan dalam sistem kontrol deteksi udara dan gas. Eksplorasi dilakukan untuk memahami prinsip kerja, karakteristik sinyal, serta cara setiap komponen berinteraksi dengan mikrokontroler. Proyek yang dibuat selanjutnya digunakan sebagai eksperimen untuk menerapkan dan membuktikan pemahaman tersebut.

\subsection{Tujuan Pembelajaran}
Setelah mempelajari laporan ini, pembaca diharapkan mampu:
\begin{enumerate}
  \item menjelaskan prinsip kerja setiap sensor, aktuator, dan komponen pendukung;
  \item menjelaskan besaran yang diukur atau dikendalikan oleh setiap komponen;
  \item menjelaskan karakteristik sinyal, pin, catu daya, dan antarmuka komponen;
  \item menjelaskan hubungan komponen dengan fitur mikrokontroler yang digunakan;
  \item menghubungkan prinsip kerja dengan rangkaian dan program sederhana;
  \item merancang eksperimen sederhana beserta rangkaiannya untuk membuktikan pemahaman terhadap komponen;
\end{enumerate}

\subsection{Ruang Lingkup}
Pembahasan mencakup komponen yang dieksplorasi oleh Focus Group, mikrokontroler
yang digunakan, rangkaian fisik, code program, serta eksperimen pembuktian.

\subsection{Daftar Komponen yang Dieksplorasi}
\begin{table}[H]
\centering
\caption{Daftar komponen yang menjadi objek eksplorasi}
\begin{tabularx}{\textwidth}{@{}clX@{}}
\toprule
No. & Komponen & Peran dalam eksplorasi \\
\midrule
1 & DHT22 & Sensor/aktuator/komponen pendukung yang dieksplorasi \\
2 & MQ-135 & Sensor/aktuator/komponen pendukung yang dieksplorasi \\
3 & MQ-2 & Sensor/aktuator/komponen pendukung yang dieksplorasi \\
4 & HC-SR501 & Sensor/aktuator/komponen pendukung yang dieksplorasi \\
5 & FC-51 & Sensor/aktuator/komponen pendukung yang dieksplorasi \\
6 & JQC-3FF-S-Z & Sensor/aktuator/komponen pendukung yang dieksplorasi \\
7 & SG90 & Sensor/aktuator/komponen pendukung yang dieksplorasi \\
8 & Micro DC Motor 6 \volt &sfsdf \\
\bottomrule
\end{tabularx}
\end{table}
